% !TEX root = ../main.tex

\chapter{نتیجه‌گیری و پیشنهادها}

% =========================================================
% 8-1
% =========================================================
\section{نتیجه‌گیری}

هدف این پروژه طراحی و پیاده‌سازی یک چت‌بات صوتی فارسی با
قابلیت اجرای محلی بود که سه مؤلفه اصلی

\begin{center}
\lr{ASR}
$\longrightarrow$
\lr{LLM}
$\longrightarrow$
\lr{TTS}
\end{center}

را در قالب یک سامانه واحد ترکیب کند.

در بخش تشخیص گفتار، چند مدل فارسی و چندزبانه از نظر دقت،
سرعت و مصرف منابع مقایسه شدند.
نتایج نشان دادند که
\lr{FastConformer Hybrid}
در مجموعه آزمایش اصلی کمترین نرخ خطای واژه را ارائه کرده و
در عین حال از سرعت مناسبی برخوردار بوده است.

در بخش مدل زبانی، چند مدل با اندازه‌های مختلف مورد بررسی
قرار گرفتند.
در ارزیابی کنترل‌شده پرسش‌های چندگزینه‌ای \lr{ParsiNLU}، مدل
\lr{Gemma 2 9B} با دقت کل \lr{46.86\%} رتبه نخست را به دست آورد،
در حالی که \lr{Qwen2.5 7B} بهترین نتیجه را در دسته ریاضی و منطق
ارائه کرد. ارزیابی درک مطلب نیز با \lr{600} نمونه \lr{ParsBench RC}
و معیارهای تطابق دقیق، \lr{F1} توکنی و شباهت معنایی طراحی شد.
آزمایش تصادفی \lr{1000} نمونه‌ای \lr{FarsInstruct} صرفاً یک مطالعه
مقدماتی بود و به‌دلیل ترکیب کنترل‌نشده دسته‌ها، مبنای رتبه‌بندی
نهایی قرار نگرفت.

در بخش تبدیل متن به گفتار نیز سه مدل فارسی مقایسه شدند.
مدل
\lr{Piper}
بهترین کیفیت ادراکی و کمترین خطای بازشناسی مجدد گفتار را
ارائه کرد، در حالی که مدل \lr{MMS} از نظر سرعت تولید بهتر بود.

بر اساس ارزیابی‌های مستقل، ترکیب زیر به‌عنوان کاندید اصلی
برای سامانه یکپارچه انتخاب شد:

\begin{center}
\lr{FastConformer Hybrid}
$\longrightarrow$
\lr{Gemma 2 2B}
$\longrightarrow$
\lr{Piper TTS}
\end{center}

این انتخاب ناظر بر محدودیت منابع اجرای محلی است و با رتبه‌بندی
کیفی مدل‌ها تعارضی ندارد؛ \lr{Gemma 2 9B} در صورت فراهم بودن منابع
بیشتر، گزینه برتر از نظر دقت چندگزینه‌ای است.

با این حال، یکی از نتایج مهم پروژه آن است که انتخاب مدل مناسب
نباید تنها بر اساس یک معیار انجام شود.
در یک سامانه گفت‌وگوی صوتی محلی، دقت، کیفیت، زمان تأخیر،
مصرف حافظه و امکان اجرای هم‌زمان مؤلفه‌ها باید به‌صورت
هم‌زمان در نظر گرفته شوند.

همچنین معماری ماژولار پروژه امکان جایگزینی هر یک از مدل‌های
\lr{ASR}، \lr{LLM} و \lr{TTS} را بدون تغییر اساسی در سایر بخش‌ها فراهم
می‌کند.
این ویژگی امکان توسعه و بهبود سامانه در آینده را ساده‌تر
خواهد کرد.


% =========================================================
% 8-2
% =========================================================
\section{محدودیت‌های پروژه}

با وجود نتایج به‌دست‌آمده، پروژه حاضر دارای چند محدودیت است.

نخست، نتایج مدل‌ها به مجموعه‌داده و شرایط آزمایش وابسته هستند
و نمی‌توان آن‌ها را به تمام سناریوهای گفتار فارسی تعمیم داد.

دوم، بخش قابل توجهی از ارزیابی کیفیت مدل زبانی و \lr{TTS} بر
معیارهای خودکار استوار بوده است.
ارزیابی انسانی گسترده می‌تواند تصویر دقیق‌تری از کیفیت واقعی
سامانه ارائه دهد.

سوم، شرایط نویزی، لهجه‌های مختلف و گفتار کاملاً محاوره‌ای
هنوز نیازمند آزمایش گسترده‌تر هستند.

چهارم، مؤلفه‌هایی مانند تشخیص فعالیت صوتی، مدیریت دقیق‌تر
تاریخچه مکالمه، پردازش نویز و ارزیابی کامل تأخیر
انتهابه‌انتها هنوز در حال تکمیل هستند.

در نتیجه، ترکیب مدل‌های انتخاب‌شده باید به‌عنوان کاندید
اصلی در وضعیت فعلی پروژه در نظر گرفته شود و نه یک انتخاب
قطعی برای تمام کاربردهای فارسی.


% =========================================================
% 8-3
% =========================================================
\section{پیشنهادها برای توسعه آینده}

برای ادامه این پروژه، چند مسیر توسعه قابل پیشنهاد است.

\begin{itemize}

    \item
    تکمیل ماژول
    \lr{VAD}
    برای تشخیص خودکار شروع و پایان گفتار؛

    \item
    بررسی روش‌های کاهش نویز و اندازه‌گیری تأثیر آن‌ها بر
    \lr{WER}
    مدل \lr{ASR}؛

    \item
    ارزیابی سامانه روی گفتار محاوره‌ای، لهجه‌های مختلف و
    محیط‌های نویزی؛

    \item
    بهبود مدیریت تاریخچه مکالمه و استفاده از روش‌هایی مانند
    خلاصه‌سازی مکالمات طولانی؛

    \item
    ارزیابی انسانی پاسخ‌های مدل زبانی و کیفیت صوت \lr{TTS}؛

    \item
    بررسی کوانتیزه‌سازی مدل زبانی برای کاهش مصرف حافظه؛

    \item
    استفاده از پردازش جریانی برای کاهش زمان شروع پاسخ؛

    \item
    اندازه‌گیری دقیق
    \lr{Time to First Token}
    و زمان شروع پخش صوت؛

    \item
    بررسی امکان شروع \lr{TTS} پیش از تکمیل کامل پاسخ مدل زبانی؛

    \item
    و توسعه یک رابط کاربری گرافیکی یا وب برای استفاده
    ساده‌تر از سامانه.

\end{itemize}


% =========================================================
% 8-4
% =========================================================
\section{ارزیابی انتهابه‌انتها در توسعه آینده}

یکی از مهم‌ترین مراحل باقی‌مانده، ارزیابی کامل سامانه در
شرایط مکالمه واقعی است.

در این ارزیابی باید زمان از پایان گفتار کاربر تا آغاز پاسخ
صوتی اندازه‌گیری شود.

به‌صورت کلی:

\begin{equation}
T_{\mathrm{E2E}}
=
T_{\mathrm{endpoint}}
+
T_{\\mathrm{ASR}}
+
T_{\\mathrm{LLM}}
+
T_{\\mathrm{TTS}}
+
T_{\mathrm{overhead}}
\end{equation}

همچنین بهتر است چندین مکالمه چندمرحله‌ای اجرا شوند تا پایداری
سامانه، مدیریت تاریخچه و انتشار خطا میان مؤلفه‌ها بررسی شود.

این آزمایش می‌تواند مشخص کند که آیا ترکیب انتخاب‌شده در
ارزیابی مستقل، در شرایط واقعی نیز بهترین گزینه باقی می‌ماند
یا نیاز به تغییر یکی از مؤلفه‌ها وجود دارد.


% =========================================================
% 8-5
% =========================================================
\section{جمع‌بندی نهایی}

در این پروژه یک چارچوب ماژولار برای توسعه چت‌بات صوتی فارسی
مبتنی بر تشخیص گفتار، مدل زبانی بزرگ و تبدیل متن به گفتار
طراحی و پیاده‌سازی شد.

مقایسه چند مدل مختلف نشان داد که انتخاب مؤلفه‌ها باید بر اساس
توازن میان کیفیت و هزینه محاسباتی انجام شود و استفاده از
بزرگ‌ترین مدل موجود الزاماً بهترین نتیجه عملی را ایجاد
نمی‌کند.

نتایج حاصل، یک ترکیب مناسب اولیه برای ساخت سامانه محلی فراهم
کرده‌اند و معماری طراحی‌شده نیز امکان ادامه توسعه، جایگزینی
مدل‌ها و افزودن قابلیت‌های جدید را فراهم می‌کند.

با تکمیل ارزیابی انتهابه‌انتها، تشخیص فعالیت صوتی، پردازش
نویز و مدیریت بهتر مکالمه، این سامانه می‌تواند به یک چت‌بات
صوتی فارسی کامل‌تر و مناسب‌تر برای تعامل بلادرنگ تبدیل شود.
